\chapter[Sensor Hardware and Integration Reference]{Sensor Hardware and Integration\\Reference}
\label{app:hardware}

A launch monitor connects sampled signals to estimates of ball and club motion, then connects those estimates to impact and flight models.
Every connection adds assumptions: the target must be identified, the measurement time and reference frame must be known, and any missing components must be supplied by additional observations or a declared model.
This appendix uses documented components as worked examples, not as an endorsed bill of materials or evidence that a particular combination achieves launch-monitor accuracy.
Read the specifications below with the measurement hierarchy in \cref{tab:hierarchy} and the impact definitions in \cref{ch:impact}.
Manufacturer specifications describe particular hardware, firmware and test conditions; the calculations here establish conditional limits, not measured performance of a completed instrument.

\section{OmniPreSense OPS243-A (24\,GHz CW Doppler)}

\subsection{Hardware and Configuration Boundaries}
The OPS243 product brief distinguishes the Doppler-only OPS243-A from the Doppler/FMCW OPS243-C: the former does not provide range from its unmodulated carrier~\cite{ops243brief}.
The brief lists a 24.00--24.25\,GHz band, 11\,dBm transmit power and a $20\degs\times24\degs$ transmit beamwidth at $-3$\,dB.
The horizontal contour width at axial distance $r$ is approximately $2r\tan(10\degs)$: 0.353\,m at 1\,m and 1.763\,m at 5\,m.
A beamwidth contour is not a hard detection boundary, and the family-level 1--100\,m motion-detection claim is not a golf-ball range guarantee.

Electrical details also require the exact variant: the USB/UART version uses 3.3\,V UART signaling, whereas an RS-232 version uses different levels.
The brief permits a 24\,V supply only for revision-C boards, otherwise 13.5\,V maximum; its 1.7\,W typical active-power table entry explicitly names OPS243-C~\cite{ops243brief}.
Do not transfer these values indiscriminately between board revisions and sensor variants or use a headline accuracy percentage as an established golf-shot uncertainty.

AN-010 documents the default 10\,ksps rate, selectable sample buffers, zero-padding, timestamped JSON and raw I/Q outputs~\cite{an010}.
For the golf example below, \texttt{S=30} sets 30\,ksps, \texttt{S(} selects 128 acquired samples and \texttt{X=32} requests a 4096-point transform; configure the buffer before the padding multiplier.
\texttt{OT}, \texttt{OM} and \texttt{OJ} select timestamps, magnitudes and JSON; \texttt{OR} supplies raw samples, while \texttt{O2} requests two reported speeds.
The UART default is 19,200 baud, 8N1.
Verify supported commands and read back the configuration for the installed firmware before saving it with \texttt{A!}.

\subsection{Sampling, Grid Spacing and Identifiability}
For an ideal monostatic CW observation, the magnitude of Doppler frequency satisfies $|f_D|=2|v_r|/\lambda$, where $v_r$ is radial velocity and $\lambda=c/f_c$.
The sign depends on the receiver convention; it must be checked before transforming observations into the book's coordinate frame.
For centered complex sampling at rate $f_s$, $N$ acquired samples and an $M$-point transform with $M\ge N$, the nominal quantities are
\begin{equation}
\label{eq:hardwaredopplergrid}
 |v_r|<\frac{\lambda f_s}{4},\qquad
 T=\frac{N}{f_s},\qquad
 \Delta v_N=\frac{\lambda f_s}{2N},\qquad
 \Delta v_M=\frac{\lambda f_s}{2M}.
\end{equation}
Here $\Delta v_N$ and $\Delta v_M$ are the unpadded and padded velocity-grid spacings.
The first expression is a sampling limit, subject also to the analog passband and firmware; it is not a guaranteed detection speed.
With $f_c=24.125$\,GHz, $c=299{,}792{,}458$\,m/s, $f_s=30{,}000$\,s$^{-1}$, $N=128$ and $M=4096$, these give 93.20\,m/s (208.48\,mph), 4.267\,ms, 1.456\,m/s (3.258\,mph) and 0.04551\,m/s (0.1018\,mph), respectively.

Zero-padding samples the same finite-record spectrum on a finer grid; it adds no observed cycles.
An isolated peak can be localized more finely than the unpadded bin spacing under suitable signal-to-noise and signal-model conditions, but neither spacing is a universal accuracy specification or a guaranteed separation between two resolvable targets.
Windowing, unequal target strengths, acceleration and clutter change the estimation problem.
This distinction explains why AN-029's nominal 0.1\,mph figure cannot establish 0.1\,mph measurement accuracy~\cite{an029}.
Likewise, a 200\,Hz reporting example has 5\,ms between outputs; that alone does not establish a 5\,ms independent observation, capture duty cycle or processing latency.

The sensor observes $v_r=\uvec{r}^{\mathsf T}\vect{v}$, where $\uvec{r}$ points from sensor to target for a positive-receding convention.
Target bearing and velocity direction are different quantities: knowing the former does not determine the latter.
If the angle $\theta$ between velocity and line of sight is independently known, $v=v_r/\cos\theta$; near transverse motion, small angular errors become strongly amplified.
For small signed perturbations, $\delta v/v\approx\delta v_r/v_r+\tan\theta\,\delta\theta$, with $\delta\theta$ in radians and nonzero denominators.
A fixed cosine correction cannot recover an unknown, changing trajectory.

\subsection{Rolling Buffer and Trigger Timing}
AN-027 describes a separate capture mode: \texttt{G1} stores 4096 complex samples in 32 segments of 128, with subsequent processing off the sensor~\cite{an027}.
A software \texttt{S!} command or 3.3\,V rising edge on J3 pin 3 (\texttt{HOST\_INT}) triggers completion.
\texttt{S\#n} selects the historical segment count; the default eight gives 1024 pre-trigger and 3072 post-trigger samples.
The record spans 136.53\,ms at 30\,ksps.
The note illustrates a microphone gate connection, but its \texttt{S\#18} example uses 250\,ksps and must not be copied as a universal golf setting.

Choose the pre-trigger history from an explicit timing budget.
For source-to-microphone distance $d$, sound speed $c_s$, trigger-chain delay $\tau_e$, desired pre-impact history $\tau_h$ and a bound $\tau_q$ on segment/timestamp uncertainty, require
\begin{equation}
\label{eq:hardwaretriggerbudget}
 \frac{128n}{f_s}\ge \frac{d}{c_s}+\tau_e+\tau_h+\tau_q,
 \qquad 0\le n\le32.
\end{equation}
At the declared illustrative value $c_s=343$\,m/s, 2\,m takes 5.831\,ms, not 6.6\,ms; the latter estimate in AN-027 is not adopted here.
Ignoring electronics and quantization only for an arithmetic example, 2\,ms of pre-impact history requires two segments at 30\,ksps, providing 8.533\,ms total history.
This is not a deployment recommendation: measure or conservatively bound the omitted delays and retain enough post-trigger data for the intended estimator.

\subsection{Interpreting the Vendor Golf Example}
AN-029 places the sensor about 2--3\,m behind the ball and gives a configuration using 30\,ksps, 128 samples, padding to 4096, mph output, a lower speed filter, a magnitude threshold and peak averaging~\cite{an029}.
Its reported 5--10\,m golf-ball detection distance is vendor application guidance, not a qualification across balls and environments.
Two reported peaks still require target association: neither peak order nor a plausible speed ratio proves which echo belongs to the ball or to a specified point on the clubhead.
The ball and rotating club can have different line-of-sight projections and different measurement times.
Before forming the smash factor of \cref{sec:smash}, establish the relevant ball/club speed definitions, timing and projection corrections.
A 1.0--1.50 acceptance gate would be an implementation heuristic, not a physical theorem; record rejected shots to avoid hiding failures through selection.
Radial-speed observations alone do not uniquely determine three-dimensional launch velocity or spin.

\section{RFbeam K-LD7 (24\,GHz FSK, Dual-RX Angle)}

Revision B describes one transmitter and two I/Q receivers separated by 6.223\,mm, with FSK range estimation and an $80\degs\times34\degs$ antenna pattern~\cite{kld7}.
Its table lists 6\,dBm typical and 10\,dBm maximum EIRP; the distance range starts at 0.005\,m, not 0.05\,m, and the nominal distance resolution is 5\,cm at the 5\,m range setting.
These are distinct entries, not demonstrated millimeter-range golf performance.
The 1$\degs$ angular-resolution entry likewise does not certify one-degree golf accuracy or full-field sensitivity.

The maximum speed setting is 100\,km/h (27.78\,m/s or 62.14\,mph), with a typical 29\,ms frame duration~\cite{kld7}.
It restricts both club and ball observations when their radial components exceed that limit; assigning the module to club speed does not remove the restriction.
More transverse placement may reduce Doppler but worsens the conversion to total speed and changes visibility.
Out-of-range targets can produce erroneous detections, rather than a clean indication that the limit was exceeded.

The host interface defaults to 115200 baud, 8E1; \texttt{INIT} can select up to 3\,Mbaud.
\texttt{GNFD} requests raw samples (\texttt{RADC}), spectra (\texttt{RFFT}), detections (\texttt{PDAT}), a tracked target (\texttt{TDAT}) and/or frame completion (\texttt{DONE})~\cite{kld7}.
Check frame counters and measured transfer time: a nominal acquisition period does not guarantee that every requested raw frame reaches the host.
Calibration must also establish antenna geometry, phase offsets, target association and timing before combining this module with another sensor.

\section{TI IWR6843 (60--64\,GHz FMCW MIMO)}

The IWR6843 has three transmitters, four receivers, a DSP and radar hardware acceleration; the datasheet distinguishes 12.5\,Msps complex-1x sampling from 25\,Msps real/complex-2x sampling~\cite{iwr6843}.
The 60--64\,GHz tuning range is not a statement that every configuration processes 4\,GHz of useful chirp bandwidth.
For an effective processed sweep bandwidth $B$, the nominal range-separation scale is $\Delta R=c/(2B)$, or 3.747\,cm at $B=4$\,GHz~\cite{tichirp}.
Windowing and signal conditions affect separation; isolated-target range precision is a different quantity.
Array geometry and the selected board determine angular coverage and aperture, so a count of twelve TX--RX channels is not twelve independent azimuth elements.

Time-division MIMO acquires the transmitter channels at different instants; velocity-induced phase must be compensated before interpreting the virtual array as an angle measurement~\cite{timimo}.
For a uniform single-target phase-sampling model with same-transmitter inter-chirp revisit interval $T_r$ within a burst, the principal unaliased radial-speed interval has half-width $\lambda/(4T_r)$, while coherent observation time sets a velocity-separation scale near $\lambda/(2T_{\rm obs})$.
Thus faster updates, wider unambiguous speed, more transmitter channels and finer Doppler separation compete for acquisition time~\cite{tichirp,timimo}.
Ambiguity-resolution methods can alter this design tradeoff but need their own assumptions and validation.

A 10--20\,Hz output stream has 50--100\,ms between reports; this is not necessarily its internal coherent dwell time.
Such reporting alone does not establish coverage of a chosen 3\,ms launch interval, and shortening the frame alone does not prove a working launch monitor.
Specify the chirps, usable bandwidth, ADC interval, transmitter schedule, synchronization, target return strength and processing/transport latency together.
Range, direction and radial velocity estimates can then constrain a trajectory model; a single range-angle-Doppler detection still lacks the tangential velocity components and does not establish spin or club delivery.

\section{Raspberry Pi Global Shutter Camera (Optical Module)}

The IMX296 camera provides 1456$\times$1088 pixels of 3.45\,\si{\micro\meter} pitch, full-resolution operation near 60\,fps and short exposures near 30\,\si{\micro\second}, subject to operating conditions~\cite{pigsproduct}.
Its XTR trigger is a \textbf{1.8\,V input}; the documented exposure is the low pulse width plus 14.26\,\si{\micro\second}.
Shared pulses can synchronize cameras, but accepted pulse rates and actual frame timing must be verified.
Raspberry Pi documents enabling the IMX296 external-trigger mode through
\begin{quote}\small\ttfamily
 echo 1 | sudo tee /sys/module/imx296/parameters/trigger\_mode
\end{quote}
Use the documented procedure for the installed driver version.
Boards with Q2 fitted require the documented R11 modification; use suitable level conversion and board-specific instructions~\cite{pigscam}.

Exposure duration controls motion blur; frame interval controls temporal sampling.
At 60\,fps, consecutive frames are 16.67\,ms apart: a ball moving at 70\,m/s travels 1.167\,m between them, but only 2.1\,mm during a 30\,\si{\micro\second} exposure.
Whether that yields usable images depends on field of view, magnification, illumination and trigger timing.
It cannot supply two consecutive frame samples within a 3\,ms interval, but it can record flight over a larger observation region.
Multiple flashes within one exposure, as discussed in \cref{app:optical}, are one possible architecture; a faster camera or a different observation window is another.
None alone establishes spin observability without resolvable surface features, timing and calibrated geometry.

\section{Simulator Integration: GSPro Open Connect}
\label{app:gspro}

GSPro Open Connect v1 is a documented JSON-over-TCP interface; its sample endpoint is localhost port 921, with the launch monitor acting as client~\cite{gspro}.
It is an integration example, not an assertion that no other documented protocol exists.
A shot supplies device identity, shot number, API version, data-presence options and ball fields including \texttt{Speed} (mph), \texttt{VLA}/\texttt{HLA} (degrees) and spin data.
The specification accepts \texttt{TotalSpin}/\texttt{SpinAxis} or \texttt{BackSpin}/\texttt{SideSpin}; it also describes optional club fields, readiness/heartbeat messages and a player-information response~\cite{gspro}.

Numeric field names do not preserve the scientific meaning of a measurement automatically.
Verify units, angle signs, handedness, reference points and event times against \cref{ch:params}; avoid labeling a signed spin-axis value as a draw independently of handedness and convention.
Two spin parameters also do not encode an unrestricted three-component spin vector; any mapping must state its representation assumptions.
The documented schema supplies no general per-field uncertainty or measurement-provenance contract.
Keep a separate, shot-linked record of raw evidence, calibration, inferred versus observed quantities and unavailable values, as motivated by \cref{ch:implications}.
Do not replace unavailable spin with an apparently measured zero or silently present a prior-driven estimate as an observation simply to fill a required field.
An accepted simulator message confirms transport compatibility, not measurement validity.

\section{Equipment Rules and Model Parameters}

Equipment conformance limits are constraints, not measured properties of every conforming club or ball.
The ball rules specify mass at most 45.93\,g and diameter at least 42.67\,mm~\cite{randaballrules}.
Using these limiting values in \cref{eq:eom} is a declared modeling choice; it does not identify a particular ball's mass, diameter, inertia or aerodynamic coefficients.

The driver spring-effect test uses a characteristic-time limit of 239\,\si{\micro\second} with an 18\,\si{\micro\second} test tolerance~\cite{usgactnotice}.
This pendulum-test quantity is not the club--ball contact duration or a universal conversion to the restitution coefficient $e$ in \cref{eq:ballspeed}.
That coefficient belongs to the declared impact model, incident speed, impact location and club/ball pairing.
The wood-head volume limit is 460\,cm$^3$ with a 10\,cm$^3$ tolerance; the MOI limit is 5900\,g\,cm$^2$ with a 100\,g\,cm$^2$ test tolerance about the vertical CG axis at a 60$\degs$ lie angle~\cite{randaclubrules}.
This is one specified inertia component, not a bound automatically transferable to every axis in \cref{eq:gear}.

\textbf{Status checked October 4, 2026:} the December 2023 announcement described changing the ODS test from 120\,mph/10$\degs$/2520\,rpm to 125\,mph/11$\degs$/2200\,rpm, originally beginning in 2028~\cite{usgaods2023}.
The June 17, 2026 joint statement superseded that timetable: no ODS testing change before January 2030 while alternative approaches are evaluated~\cite{usgaods2026}.
Do not present the earlier 2028 schedule as current or assume the eventual conditions are settled.
The 317\,yd limit plus 3\,yd tolerance concerns the specified overall-distance test, not a universal maximum carry for a conforming ball~\cite{randaballrules}.
A flight-model comparison requires the actual launch conditions, atmosphere, aerodynamic characterization and prescribed bounce/roll treatment; merely keeping an arbitrary simulation below 320\,yd neither validates it nor certifies ball conformance.

\section{Detection Theory: CFAR Selection}

Constant-false-alarm-rate (CFAR) detection adapts a threshold to an estimate of local background power.
A fixed SNR cutoff and a masked region around zero Doppler do not, by themselves, specify a cell-averaging CFAR detector.
CA-CFAR uses the mean of training-cell powers, excluding guard cells around the cell under test; OS-CFAR uses a selected order statistic~\cite{cfartutorial,cfarmethods}.
The training region, guard width, rank where applicable, noise model and target false-alarm probability are part of the algorithm.

For an illustrative square-law CA detector, suppose the test-cell noise and $K$ training-cell powers are independent exponential variables with common mean $\mu$.
Let $\bar P$ be their training mean and declare a detection when $P_{\rm CUT}>\alpha\bar P$.
Conditioning on the training samples gives
\begin{equation}
\label{eq:hardwarecfar}
 P_{\rm FA}=\mathbb E\!\left[e^{-\alpha\bar P/\mu}\right]
 =\left(1+\frac{\alpha}{K}\right)^{-K},\qquad
 \alpha=K\left(P_{\rm FA}^{-1/K}-1\right).
\end{equation}
The product follows from the exponential Laplace transform for each independent training cell.
It is a conditional model calculation, not a false-alarm guarantee for real golf data: windowing and zero-padding can correlate spectral samples, while clutter and target leakage can invalidate the common-background model.
A mask specified only as a number of padded bins also changes its physical velocity width when the FFT configuration changes.

OS-CFAR can reduce the effect of some high-power contaminated training cells, with performance depending on the chosen rank and contamination pattern; it is not universally superior for a club followed by a ball~\cite{cfarmethods}.
Their closeness in physical space or time does not establish closeness in the processed range--Doppler cells.
Guard cells cannot separate unresolved targets, and neither a universal 0.5--1\,dB loss nor a preferred detector follows from the shot geometry alone.
Compare candidate detectors at matched false-alarm targets on declared single-target, multiple-target and clutter cases, including nondetections and mistaken associations.

\begin{implication}
The useful design unit is the complete observation chain: sensor geometry, acquisition window, synchronization, estimator, uncertainty record and downstream model.
A finer reported number does not create a new physical observation; a plausible impact or flight result does not independently validate the measurements used to produce it.
Preserve these distinctions when relating measured ball flight back to club delivery and, further upstream, to the golfer's swing.
\end{implication}
